\cotefile{57}
\scanpage{133}
% archive foliation: 132
\margin{$\ast$}
\uline{Théorème} Soit \deleted{$\unsure{Y'}\longleftarrow\unsure{X'}$}\inserted{$f:X\to Y$} un morphisme \uline{surjectif} de préschémas propres sur un corps $k$. Alors le noyau $\uuline{N}$ de $\uuline{\mathrm{Pic}}_{\unsure{Y}/k}\longrightarrow\uuline{\mathrm{Pic}}_{X/k}$ est \deleted{contenu} un sous-groupe d'un \uline{groupe affine} $G$ sur $k$.
\margin{?}

\unsure{Donc}

\uline{Corollaire}. Si $\uline{\mathrm{Pic}}_{Y/k}$ et $\uline{\mathrm{Pic}}_{X/k}$ existent, le noyau de $\uline{\mathrm{Pic}}_{Y/k}\mathbin{\text{---}}\uline{\mathrm{Pic}}_{X/k}$ est un groupe alg.-affine.

\uline{Corollaire} Dans les conditions précédentes, $\unsure{\mathrm{Nér\,Sév}}_{0}(Y/\deleted{\unsure{k}}k)\mathbin{\text{---}}\unsure{\mathrm{Nér\,Sév}}_{0}(X/k)$ est \uline{injectif}, et si $A$ et $B$ sont les ``parties abéliennes'' de $\uline{\mathrm{Pic}}_{Y/k}$ et $\uline{\mathrm{Pic}}_{X/k}$ \inserted{[\unsure{définies} si $k$ est parfait]} alors $A\longrightarrow B$ a un noyau fini.

\uline{Dém} \deleted{$\uline{\mathrm{Pic}}_X\unsure{\longleftarrow}\uline{\mathrm{Pic}}_{\unsure{Y}}$}, On va pour simplifier admettre que tous les $\uuline{\mathrm{Pic}}$ envisagés sont représentables.

a) On peut supposer $k$ \uline{alg. clos} (\unsure{extension}? )

b) On aura un diagramme de \uuline{schémas} \unsure{sur $k$}
\diagram{\left.\begin{tikzcd}
\uline{\mathrm{Pic}}_{X/k} \arrow[d,dash] & \uline{\mathrm{Pic}}_{Y/\unsure{k}} \arrow[l] \arrow[d,"\uline{\text{Aff.}}"] \\
\uline{\mathrm{Pic}}_{X_{\mathrm{réd}}} & \uline{\mathrm{Pic}}_{Y_{\mathrm{réd}}} \arrow[l,"?"]
\end{tikzcd}\right]\mathbin{\text{---}}\unsure{\text{ont}}}
